\documentclass[11pt,a4paper]{article}

\usepackage[utf8]{inputenc}
\usepackage[T1]{fontenc}
\usepackage[spanish,es-noshorthands]{babel}
\usepackage{lmodern}
\usepackage[margin=2.3cm]{geometry}
\usepackage{booktabs}
\usepackage{tabularx}
\usepackage{array}
\usepackage{xcolor}
\usepackage{enumitem}
\usepackage{parskip}
\usepackage{listings}
\usepackage{titlesec}
\usepackage{fancyhdr}
\usepackage[colorlinks=true,linkcolor=azul,urlcolor=azul,citecolor=azul]{hyperref}

% ---- Colores ----
\definecolor{azul}{HTML}{1d4ed8}
\definecolor{indigo}{HTML}{4f46e5}
\definecolor{gris}{HTML}{f3f4f6}
\definecolor{gristxt}{HTML}{4b5563}

% ---- Encabezados ----
\titleformat{\section}{\Large\bfseries\color{indigo}}{\thesection.}{0.5em}{}
\titleformat{\subsection}{\large\bfseries\color{azul}}{\thesubsection}{0.5em}{}

% ---- Pie de página ----
\pagestyle{fancy}
\fancyhf{}
\renewcommand{\headrulewidth}{0pt}
\fancyfoot[L]{\small\color{gristxt}Traductor de audio en tiempo real}
\fancyfoot[R]{\small\color{gristxt}\thepage}

% ---- Estilo de bloques de código ----
\lstdefinestyle{cmd}{
  basicstyle=\ttfamily\small,
  backgroundcolor=\color{gris},
  frame=tb,
  framesep=5pt,
  framerule=0pt,
  aboveskip=8pt, belowskip=8pt,
  breaklines=true,
  columns=fullflexible,
  keepspaces=true,
  showstringspaces=false,
}
\lstset{style=cmd}

% Columna tipo X centrada y X angosta para tablas
\newcolumntype{Y}{>{\raggedright\arraybackslash}X}

\hypersetup{pdftitle={Manual de usuario - Traductor de audio en tiempo real}}

\begin{document}

% ======== PORTADA ========
\begin{center}
  {\color{indigo}\rule{\textwidth}{2pt}}\\[1.2em]
  {\Huge\bfseries Traductor de audio en tiempo real}\\[0.6em]
  {\Large\color{gristxt} Manual de usuario}\\[1em]
  {\large Traduce lo que suena en tu PC a otro idioma, en vivo,\\
   con \textbf{Gemini 3.5 Live Translate}}\\[1em]
  {\color{indigo}\rule{\textwidth}{2pt}}
\end{center}

\vspace{1em}
\noindent\textit{Ejemplo: un curso de Udemy en inglés lo escuchas en español
mientras lo ves, con subtítulos en pantalla.}

\vspace{1.5em}
\tableofcontents
\vspace{1em}
\hrule
\vspace{1em}

% ======== 1. REQUISITOS ========
\section{¿Qué necesito? (requisitos)}

\begin{tabularx}{\textwidth}{@{}l Y l@{}}
\toprule
\textbf{Necesitas} & \textbf{Para qué} & \textbf{Costo} \\
\midrule
Windows 10 u 11 & Sistema operativo & --- \\
Python 3.10 o superior & Ejecutar el programa & Gratis \\
VB-CABLE & Cable de audio virtual para capturar el sonido & Gratis \\
Clave de Gemini (API key) & Conectar con el traductor de Google & Gratis (tier de prueba) \\
Audífonos o bocinas & Escuchar la traducción & --- \\
Internet & La traducción ocurre en la nube & --- \\
\bottomrule
\end{tabularx}

% ======== 2. INSTALACIÓN ========
\section{Instalación paso a paso}

\subsection{Instalar Python}
\begin{enumerate}[leftmargin=*]
  \item Descarga Python desde \url{https://www.python.org/downloads/}
  \item Ejecuta el instalador y \textbf{marca la casilla ``Add Python to PATH''} (importante).
  \item Clic en \textbf{Install Now}.
\end{enumerate}
Para verificar, abre \textbf{PowerShell} y escribe:
\begin{lstlisting}
python --version
\end{lstlisting}
Debe mostrar algo como \texttt{Python 3.12.x}.

\subsection{Instalar VB-CABLE (el cable de audio)}
\begin{enumerate}[leftmargin=*]
  \item Ve a \url{https://vb-audio.com/Cable} y descarga el paquete.
  \item Descomprime el ZIP.
  \item \textbf{Clic derecho} en \texttt{VBCABLE\_Setup\_x64.exe} $\rightarrow$ \textbf{Ejecutar como administrador}.
  \item Clic en \textbf{Install Driver} y \textbf{reinicia} la computadora.
\end{enumerate}
Esto agrega dos dispositivos: \textbf{CABLE Input} (entrada) y \textbf{CABLE Output}
(salida). Es seguro y se desinstala con el mismo instalador (\textit{Remove Driver}).

\subsection{Descargar el programa}
\textbf{Opción fácil (sin Git):} entra al repositorio, botón verde \textbf{Code}
$\rightarrow$ \textbf{Download ZIP} y descomprime donde quieras.

\noindent\textbf{Opción con Git:}
\begin{lstlisting}
git clone https://github.com/joseIronMaker/traductor-audio-tiempo-real.git
\end{lstlisting}

\subsection{Instalar las dependencias}
Abre \textbf{PowerShell dentro de la carpeta del programa} (clic derecho en la
carpeta $\rightarrow$ \textit{Abrir en Terminal}) y ejecuta:
\begin{lstlisting}
python -m venv .venv
.\.venv\Scripts\Activate.ps1
pip install -r requirements.txt
\end{lstlisting}
Si al activar sale un error de ``ejecución de scripts deshabilitada'', corre esto
una vez y vuelve a intentar:
\begin{lstlisting}
Set-ExecutionPolicy -Scope CurrentUser RemoteSigned
\end{lstlisting}

\subsection{Conseguir tu clave de Gemini}
\begin{enumerate}[leftmargin=*]
  \item Entra a \url{https://aistudio.google.com/apikey} (inicia sesión con Google).
  \item Clic en \textbf{Create API key} y \textbf{copia} la clave.
  \item Copia el archivo \texttt{.env.example} y renómbralo a \texttt{.env}:
  \begin{lstlisting}
copy .env.example .env
  \end{lstlisting}
  \item Abre \texttt{.env} con el Bloc de notas y pega tu clave:
  \begin{lstlisting}
GEMINI_API_KEY=pega_aqui_tu_clave
  \end{lstlisting}
  \item Guarda. \textbf{Nunca compartas este archivo} (contiene tu clave).
\end{enumerate}

\subsection{Probar que todo conecta}
\begin{lstlisting}
python hola_mundo.py
\end{lstlisting}
Si dice \textbf{``Conexión con la Live API establecida''}, todo quedó listo.

% ======== 3. AUDIO ========
\section{Configurar el audio (lo más importante)}

El truco es mandar \textbf{el sonido que quieres traducir} al cable
(CABLE Input). El programa lo captura de \textbf{CABLE Output}, lo traduce, y
reproduce el resultado en tus audífonos.

\begin{lstlisting}
Fuente (curso/reunion) -> CABLE Input -> [cable] -> CABLE Output -> el programa
                                                                       |
                                                                    traduce
                                                                       v
            oyes la traduccion  <-  tus audifonos/bocinas  <----------
\end{lstlisting}

\noindent\textbf{Regla de oro:} la \textbf{fuente} va a \textit{CABLE Input};
tú \textbf{escuchas} por tus audífonos. La ``Salida de audio'' del programa
\textbf{nunca} debe ser un cable.

\subsection{Cómo enrutar según lo que uses}
\textbf{A) Un navegador (YouTube, Udemy, Coursera):}
\begin{enumerate}[leftmargin=*,nosep]
  \item Clic derecho en el ícono de volumen (barra de tareas) $\rightarrow$ \textbf{Mezclador de volumen}.
  \item Busca tu navegador $\rightarrow$ \textbf{Dispositivo de salida} = \textbf{CABLE Input}.
\end{enumerate}

\textbf{B) Zoom:}
\begin{enumerate}[leftmargin=*,nosep]
  \item En Zoom, flecha junto al micrófono $\rightarrow$ \textbf{Seleccionar un altavoz}.
  \item Elige \textbf{CABLE Input (VB-Audio Virtual Cable)}.
\end{enumerate}

\textbf{C) Todo el sistema:}
\begin{enumerate}[leftmargin=*,nosep]
  \item Clic derecho en volumen $\rightarrow$ \textbf{Configuración de sonido}.
  \item En \textbf{Salida}, elige \textbf{CABLE Input} como predeterminado.
  (Recuerda regresarlo a tus audífonos al terminar.)
\end{enumerate}

% ======== 4. USO ========
\section{Usar la aplicación}

Con el entorno activado, abre la app:
\begin{lstlisting}
python app.py
\end{lstlisting}

Configura la ventana así:

\begin{tabularx}{\textwidth}{@{}l Y@{}}
\toprule
\textbf{Campo} & \textbf{Qué poner} \\
\midrule
Idioma origen & Déjalo en \texttt{Auto-detectar} (el modelo lo detecta solo). \\
Idioma destino & El idioma al que quieres traducir (ej. \texttt{Español}). \\
Salida de audio & Tus audífonos o bocinas (donde escuchas). \\
Capturar de & \texttt{CABLE Output (VB-Audio Virtual Cable)} (por defecto). \\
\bottomrule
\end{tabularx}

\vspace{0.8em}
\noindent\textbf{Botones:}
\begin{itemize}[leftmargin=*,nosep]
  \item \textbf{Iniciar} --- empieza a traducir (estado \textit{Traduciendo}, en verde).
  \item \textbf{Pausar / Reanudar} --- detiene/continúa (ahorra costo en pausa).
  \item \textbf{Detener} --- termina la sesión.
  \item \textbf{Limpiar} --- borra los subtítulos en pantalla.
  \item \textbf{Mostrar inglés} --- muestra también el texto del idioma original.
\end{itemize}

\vspace{0.5em}
\noindent\textbf{Cambiar de idioma en vivo:} mientras corre, solo cambia el menú
\textbf{Idioma destino} y en un par de segundos se adapta solo, sin reiniciar.

Cuando pongas a sonar la fuente, en unos segundos verás los subtítulos avanzar y
oirás la traducción.

% ======== 5. CASOS DE USO ========
\section{Casos de uso}

\begin{tabularx}{\textwidth}{@{}l Y@{}}
\toprule
\textbf{Quiero traducir...} & \textbf{Pasos} \\
\midrule
Un video de YouTube & Rutea el navegador a CABLE Input, destino = tu idioma, Iniciar. \\
Un curso de Udemy/Coursera & Igual que YouTube (es el navegador). \\
Una reunión de Zoom & Zoom: Speaker = CABLE Input, Iniciar. \\
Cualquier cosa del sistema & Salida predeterminada = CABLE Input. \\
\bottomrule
\end{tabularx}

% ======== 6. PROBLEMAS ========
\section{Problemas comunes y soluciones}

\begin{tabularx}{\textwidth}{@{}l Y Y@{}}
\toprule
\textbf{Síntoma} & \textbf{Causa probable} & \textbf{Solución} \\
\midrule
Error \texttt{-9999} al Iniciar & Dispositivo WASAPI/WDM-KS conflictivo (USB) & La app usa MME; elige otra Salida (ej. Speakers). No uses dispositivos ``WDM-KS''. \\
\addlinespace
Error \texttt{-9985} (Device unavailable) & Otra app (ej. Zoom) tiene tomado el dispositivo & Usa otra salida o cierra/redirige la otra app. \\
\addlinespace
La traducción se repite o se atora & Hay un eco en el ruteo & \texttt{Win+R} $\rightarrow$ \texttt{mmsys.cpl} $\rightarrow$ Grabar $\rightarrow$ CABLE Output $\rightarrow$ Escuchar $\rightarrow$ desmarca ``Escuchar este dispositivo''. \\
\addlinespace
No escucho la traducción & La salida está en un cable o en un dispositivo que no usas & Pon Salida de audio = tus audífonos/bocinas reales. \\
\addlinespace
No aparecen subtítulos & La fuente no llega al cable & Revisa que enrutaste el navegador/Zoom a CABLE Input. \\
\addlinespace
``Falta GEMINI\_API\_KEY'' & No creaste el \texttt{.env} o está vacío & Repite la sección 2.5. \\
\addlinespace
Va con retraso de unos segundos & Normal & Es la latencia del modelo en vivo; en clases grabadas no afecta. \\
\bottomrule
\end{tabularx}

% ======== 7. FAQ ========
\section{Preguntas frecuentes}

\begin{description}[leftmargin=0pt,style=nextline]
  \item[¿Es gratis?] El programa sí. Gemini tiene un tier gratuito para probar;
  en plan de pago cuesta aprox. \$0.037 USD por minuto de audio (aprox.
  \$2.2/hora). Pausar cuando no necesitas traducir ahorra mucho.
  \item[¿Oiré también el audio original?] No: el original se va ``por dentro''
  al cable y solo oyes la traducción (así no se encima).
  \item[¿A qué idiomas traduce?] A más de 70. El menú incluye español, inglés,
  francés, alemán, portugués, italiano, japonés, coreano, chino, ruso, árabe e hindi.
  \item[¿Detecta solo el idioma original?] Sí, automáticamente. Por eso
  ``Idioma origen'' es solo informativo.
  \item[¿Funciona en Mac o Linux?] El código de Gemini sí, pero la captura está
  pensada para Windows (VB-CABLE). En Mac se usa BlackHole y en Linux el monitor
  de PulseAudio.
  \item[¿Es seguro VB-CABLE?] Sí, es software conocido de VB-Audio. Descárgalo
  solo del sitio oficial \texttt{vb-audio.com}.
\end{description}

\vspace{1em}
\hrule
\vspace{0.6em}
\noindent\small\color{gristxt} Repositorio:
\url{https://github.com/joseIronMaker/traductor-audio-tiempo-real}

\end{document}
